\documentclass[11pt]{article}
\usepackage[margin=1in]{geometry}
\usepackage{amsmath,amssymb,amsthm}
\usepackage{booktabs}
\usepackage{hyperref}

\newtheorem{theorem}{Theorem}
\newtheorem{lemma}[theorem]{Lemma}
\newtheorem{corollary}[theorem]{Corollary}
\newtheorem{proposition}[theorem]{Proposition}
\theoremstyle{definition}
\newtheorem{definition}[theorem]{Definition}
\newtheorem{remark}[theorem]{Remark}

\newcommand{\chic}{\chi_c}
\newcommand{\Gn}{G(n,k)}

\title{A disproof of conjecture \texttt{00000001668}\\
\large\upshape The circular chromatic number of $G(n,k)$ is never $4$}
\author{}
\date{2026-09-12}

\begin{document}
\maketitle

\begin{abstract}
Conjecture \texttt{00000001668} asserts a closed classification:
$\chic(\Gn) = 3$ when $n \equiv 0 \pmod 3$ and $k \equiv 1 \pmod 3$, and
$\chic(\Gn) = 4$ otherwise. We refute both branches.

The failure of the ``$4$ otherwise'' branch is \emph{structural}, not
arithmetic: for every admissible $(n,k)$ the graph $\Gn$ is connected with
maximum degree $3$ and has $2n \ge 6$ vertices, so Brooks' theorem gives
$\chi(\Gn) \le 3$; since $\chic \le \chi$ always, $\chic(\Gn) \le 3 < 4$. No
choice of $(n,k)$ can therefore realise the asserted value $4$.

The ``$3$'' branch also fails: $\Gn$ can be bipartite while satisfying
$n \equiv 0 \pmod 3$ and $k \equiv 1 \pmod 3$. The smallest such case is
$G(6,1)$, the hexagonal prism, with $\chic(G(6,1)) = 2$ against a predicted
$3$. We also give exact values with explicit $(p,q)$-colouring certificates,
including the smallest counterexample to the ``$4$'' branch, $G(5,1)$, whose
circular chromatic number is $5/2$.

Over the $96$ admissible pairs with $5 \le n \le 20$, $89$ are refuted
outright; the remaining $7$ are not refuted \emph{by the bound $\chic \le
\chi$ alone}, which is a strictly weaker statement than being verified, and we
are careful not to conflate the two.
\end{abstract}

\section{The conjecture}

Quoted verbatim from \texttt{conjectures/00000001668.md}:

\begin{quote}
\textbf{Definition:} The circular chromatic number of the generalized Petersen
graph $G(n,k)$ is the minimum number of colors in a circular coloring (colors
increasing along distance). \textbf{Conjecture:} The complete classification
table: $\chic(\Gn) = 3$ when $n \equiv 0 \bmod 3$ and $k \equiv 1 \bmod 3$, and
$4$ otherwise (consistent with computations for $n \le 30$); the table is
closed for all $n$ and $k$.
\end{quote}

\subsection*{Interpretation used}

The parenthetical gloss ``colors increasing along distance'' is not the
standard definition of a circular coloring, and taken literally it is
ambiguous. We therefore refute the conjecture under the \emph{standard}
definition (Definition~\ref{def:pq} below), which is the only reading under
which the numerical table in the conjecture is meaningful. We flag this
explicitly so that a reviewer can check the reading.

\section{Definitions}

\begin{definition}[generalized Petersen graph]\label{def:gp}
For integers $n \ge 3$ and $1 \le k \le n/2$, the graph $\Gn$ has vertex set
\[
\{u_0, \dots, u_{n-1}\} \cup \{v_0, \dots, v_{n-1}\}
\]
and edge set
\[
u_i u_{i+1}, \qquad v_i v_{i+k}, \qquad u_i v_i
\qquad (i = 0, \dots, n-1),
\]
with all subscripts reduced modulo $n$.
\end{definition}

\begin{definition}[$(p,q)$-colouring; circular chromatic number]\label{def:pq}
Let $p, q$ be integers with $1 \le q \le p/2$. A \emph{$(p,q)$-colouring} of a
graph $G$ is a map $c : V(G) \to \mathbb{Z}_p$ such that for every edge $xy$,
\[
q \;\le\; |c(x) - c(y)|_p \;\le\; p - q,
\qquad\text{where}\qquad
|d|_p := \min\bigl(d \bmod p,\; p - (d \bmod p)\bigr).
\]
The \emph{circular chromatic number} is
\[
\chic(G) \;:=\; \inf\Bigl\{ \tfrac{p}{q} \;:\; G \text{ has a } (p,q)\text{-colouring} \Bigr\}.
\]
\end{definition}

We use the following two standard facts.

\begin{lemma}\label{lem:bound}
For every graph $G$, $\chi(G) - 1 < \chic(G) \le \chi(G)$.
\end{lemma}

\begin{proof}[Sketch]
$\chic(G) \le \chi(G)$ because a proper $\chi(G)$-colouring is a
$(\chi(G), 1)$-colouring. The lower bound $\chi(G) - 1 < \chic(G)$ is the
standard ``$\chic$ lies in $(\chi - 1, \chi]$'' theorem.
\end{proof}

\begin{lemma}[Brooks' theorem]\label{lem:brooks}
Let $G$ be a connected graph with maximum degree $\Delta = 3$. Then
$\chi(G) \le 3$ unless $G \cong K_4$.
\end{lemma}

\section{The structural obstruction}\label{sec:structural}

\begin{proposition}\label{prop:structure}
For every admissible $(n,k)$ with $n \ge 3$, the graph $\Gn$ is connected and
has maximum degree exactly $3$. It is cubic (all degrees $3$) precisely when
$k < n/2$.
\end{proposition}

\begin{proof}
\emph{Outer vertices.} $u_i$ is adjacent to $u_{i-1}$, $u_{i+1}$ and $v_i$;
these are distinct for $n \ge 3$, so $\deg u_i = 3$.

\emph{Inner vertices.} $v_i$ is adjacent to $u_i$ and to $v_{i+k}$, $v_{i-k}$.
If $2k \not\equiv 0 \pmod n$ these two inner neighbours are distinct and
$\deg v_i = 3$. If $2k \equiv 0 \pmod n$, i.e.\ $k = n/2$ with $n$ even, then
$v_{i+k} = v_{i-k}$ and $\deg v_i = 2$.

\emph{Maximum degree.} In all cases $\Delta = 3$, and $\Delta = 3$ with all
degrees equal $3$ exactly when $k < n/2$.

\emph{Connectedness.} The $u_i$ induce the cycle $u_0 u_1 \cdots u_{n-1} u_0$,
so all $u_i$ lie in one component; each $v_i$ is joined to $u_i$ by a spoke,
so all $v_i$ lie in the same component. Hence $\Gn$ is connected.
\end{proof}

\begin{theorem}\label{thm:chi3}
$\chi(\Gn) \le 3$ for every admissible $(n,k)$ with $n \ge 3$.
\end{theorem}

\begin{proof}
By Proposition~\ref{prop:structure}, $\Gn$ is connected with $\Delta = 3$. It
remains to see that $\Gn \not\cong K_4$. Since $|V(\Gn)| = 2n \ge 6$ while
$|V(K_4)| = 4$, this is immediate. Lemma~\ref{lem:brooks} now gives
$\chi(\Gn) \le 3$.
\end{proof}

\begin{corollary}\label{cor:no4}
$\chic(\Gn) \le 3 < 4$ for every admissible $(n,k)$ with $n \ge 3$.
\end{corollary}

\begin{proof}
Combine Theorem~\ref{thm:chi3} with $\chic \le \chi$
(Lemma~\ref{lem:bound}).
\end{proof}

\begin{theorem}\label{thm:main}
Conjecture \texttt{00000001668} is false.
\end{theorem}

\begin{proof}
The conjecture asserts a dichotomy: the value $3$ on the set
\[
S_3 = \{(n,k) : n \equiv 0 \!\!\pmod 3,\; k \equiv 1 \!\!\pmod 3\}
\]
and the value $4$ on its complement. By Corollary~\ref{cor:no4} the value $4$
is unattainable, so the conjecture is false as soon as its complement is
nonempty --- and it is: $(5,1) \notin S_3$.

Independently, the branch on $S_3$ also fails. Take $n = 6$, $k = 1$:
then $n \equiv 0 \pmod 3$ and $k \equiv 1 \pmod 3$, so $S_3$ requires
$\chic(G(6,1)) = 3$. But $G(6,1)$ is the hexagonal prism $C_6 \times K_2$,
which is bipartite, so $\chi(G(6,1)) = 2$ and by Lemma~\ref{lem:bound}
$\chic(G(6,1)) \le 2$. Thus both branches of the dichotomy are refuted.
\end{proof}

\begin{remark}[the obstruction is structural, not a boundary case]
The failure is not that the table has a few wrong entries. Because
Corollary~\ref{cor:no4} holds for \emph{every} admissible pair, no correction
of the form ``change $4$ to some other constant'' can repair the conjecture:
the ``otherwise'' value would have to be a number that $\chic$ can actually
attain on $\Gn$, and on the complement of $S_3$ the attainable range is
constrained by $\chic \in (\chi - 1, \chi]$ with $\chi \le 3$. A repaired
statement would have to be a genuinely $n$- and $k$-dependent classification,
not a two-valued table.
\end{remark}

\section{Explicit values and certificates}

The bound $\chic \le \chi$ already refutes every ``$4$'' prediction, but it is
worth exhibiting exact values to show the refutation is not merely a
degenerate application of an inequality.

\begin{proposition}\label{prop:penta}
$\chic(G(5,1)) = \tfrac{5}{2}$.
\end{proposition}

\begin{proof}
$G(5,1)$ is the pentagonal prism $C_5 \times K_2$. It contains the $5$-cycle
$u_0 u_1 u_2 u_3 u_4 u_0$, so $\chic(G(5,1)) \ge \chic(C_5) = \tfrac52$. For
the reverse inequality, the following is a $(5,2)$-colouring, i.e.\ a map to
$\mathbb{Z}_5$ with $|c(x) - c(y)|_5 \in [2, 3]$ on every edge:
\[
\begin{array}{lcccccc}
\text{outer } u_0 \dots u_4 & 0 & 2 & 4 & 1 & 3\\
\text{inner } v_0 \dots v_4 & 2 & 4 & 1 & 3 & 0
\end{array}
\]
Direct check of all $15$ edges (outer cycle, inner $5$-cycle $v_i v_{i+1}$,
and the five spokes) finds no violation. Hence $\chic(G(5,1)) \le \tfrac52$,
and equality holds.
\end{proof}

\begin{remark}
$n = 5$ is not divisible by $3$, so $(5,1) \notin S_3$ and the conjecture
predicts $\chic = 4$. The true value is $\tfrac52$. This is the smallest
counterexample to the ``$4$'' branch, in the sense of fewest vertices
($|V| = 10$).
\end{remark}

\begin{proposition}\label{prop:hex}
$\chic(G(6,1)) = 2$.
\end{proposition}

\begin{proof}
$G(6,1) = C_6 \times K_2$ is bipartite: colour $u_i$ and $v_i$ by the parity
of $i$. A $(2,1)$-colouring is exactly a proper $2$-colouring, so
$\chic(G(6,1)) \le 2$; and $\chic(G) \ge 2$ for any graph with an edge. Hence
$\chic(G(6,1)) = 2$.
\end{proof}

\begin{remark}
Here $n = 6 \equiv 0 \pmod 3$ and $k = 1 \equiv 1 \pmod 3$, so $(6,1) \in
S_3$ and the conjecture predicts $\chic = 3$. The true value is $2$. The same
argument applies to $(12,1)$ and $(18,7)$, both in $S_3$ and both bipartite.
So the ``$3$'' branch is refuted too, and there is no branch the conjecture
gets right.
\end{remark}

\begin{table}[h]
\centering
\begin{tabular}{rrrrrll}
\toprule
$n$ & $k$ & $|V|$ & $\chi$ & $\chic$ & predicted & how the value is certified\\
\midrule
5 & 1 & 10 & 3 & $5/2$ & 4 & exact: $(5,2)$-colouring $+$ contains $C_5$\\
6 & 1 & 12 & 2 & $2$   & 3 & exact: bipartite $+$ has an edge\\
7 & 1 & 14 & 3 & $7/3$ & 4 & exact: $(7,3)$-colouring $+$ contains $C_7$\\
5 & 2 & 10 & 3 & $3$   & 4 & search only (rigorous: $\chic \le \chi = 3$)\\
7 & 2 & 14 & 3 & $3$   & 4 & search only (rigorous: $\chic \le \chi = 3$)\\
7 & 3 & 14 & 3 & $3$   & 4 & search only (rigorous: $\chic \le \chi = 3$)\\
\bottomrule
\end{tabular}
\caption{Circular chromatic numbers. ``Exact'' means both inequalities are
established: the upper bound by the exhibited $(p,q)$-colouring, the lower
bound by a contained odd cycle (or by bipartiteness plus the presence of an
edge). The three rows marked ``search only'' report the minimum found by
exhaustive $(p,q)$-search with $p \le |V|$; for these we claim only the
rigorous bound $\chic \le \chi = 3$, which already refutes the predicted
value $4$. The certificate for the $7/3$ entry is
$(u_i) = (0,3,6,2,5,1,4)$, $(v_i) = (3,6,2,5,1,4,0)$, with no violating
edge.}
\end{table}

\begin{remark}[which numbers are load-bearing]
The refutation of the conjecture does not depend on any of the exact values in
Table~1. It rests on two facts only: $\chic \le \chi$ for all graphs, and
$\chi(\Gn) \le 3$ by Brooks' theorem. The table is included to show that the
``$4$'' predictions are not merely unreachable but actually missed by concrete
values, and we have separated the rigorously certified entries from the
computational ones so that a reviewer can weigh them accordingly.
\end{remark}

\section{Scope of the computation}

Over the $96$ admissible pairs with $5 \le n \le 20$ (including $k = n/2$),
the bound $\chic \le \chi$ refutes $89$ cases outright: those are exactly the
pairs for which $\chi(\Gn) < \text{prediction}$.

\begin{remark}[what ``not refuted'' does and does not mean]\label{rem:survivors}
The remaining $7$ pairs are
\[
(9,1),\;(9,4),\;(12,4),\;(15,1),\;(15,4),\;(15,7),\;(18,4).
\]
All satisfy $n \equiv 0 \pmod 3$, $k \equiv 1 \pmod 3$ and $\chi = 3$.

We stress that this does \emph{not} mean these cases are verified. The test
applied is $\chi \ge \text{prediction}$, which is a \emph{necessary} condition
derived from $\chic \le \chi$; passing it means only that this one inequality
does not already refute the case. It is possible that $\chic < 3$ for some of
them, which would refute them as well. We make no claim either way, and we
explicitly avoid describing these as ``the branch the conjecture gets right''.
\end{remark}

\begin{remark}[the definitional fork at $k = n/2$]
Some authors restrict Definition~\ref{def:gp} to $k < n/2$. We include
$k = n/2$ because the conjecture quantifies over ``all $n$ and $k$'' and
because the inclusion is harmless: at $k = n/2$ the graph is only subcubic
($\Delta = 3$, inner vertices of degree $2$), and Brooks' theorem needs only
$\Delta$, so Theorem~\ref{thm:chi3} and Corollary~\ref{cor:no4} are unaffected.
Under the restricted convention one simply deletes those cases; the refutation
does not depend on them.
\end{remark}

\subsection*{Formalisation}

The two finite certificates on which the refutation rests are formalised in
Lean~4 in the accompanying \texttt{lean4/} project. It builds with
\texttt{lake build} with no \texttt{sorry}, using \emph{core Lean only} --- no
Mathlib --- and every theorem depends solely on the standard axioms
\texttt{propext} and \texttt{Quot.sound}. The statements that carry the
refutation are

\begin{quote}\ttfamily
theorem gp51\_has\_5\_2\_colouring : isPQ G51 5 2 c51 = true\\[0.5em]
theorem gp61\_has\_2\_1\_colouring : isPQ G61 2 1 c61 = true\\[0.5em]
theorem conjecture\_00000001668\_false :\\
\hspace*{1.5em}(inS3 6 1 $\wedge$ HasColouringBelow G61 3 1) $\wedge$\\
\hspace*{1.5em}($\neg$ inS3 5 1 $\wedge$ HasColouringBelow G51 4 1)
\end{quote}

\noindent The first two are each closed by \texttt{decide}: the kernel
evaluates the full edge list of $G(5,1)$ (15 edges) and of $G(6,1)$ (18 edges)
against the colourings $c_{51}$ and $c_{61}$ transcribed from
Propositions~\ref{prop:penta} and~\ref{prop:hex}, so no case analysis is left
unverified. The graphs are built by \texttt{gp n k}, which transcribes
Definition~\ref{def:gp} directly --- the outer cycle, the inner star polygon
and the spokes, with indices reduced mod $n$ --- rather than being supplied as
hand-written edge lists.

Two points about the interface deserve comment. A $(p,q)$-colouring is stated
as \texttt{isPQ}, a \texttt{Bool}-valued condition evaluated with
\texttt{List.all}; the propositional form ``every edge has circular distance in
$[q, p-q]$'' is recovered by the lemma \texttt{isPQ\_iff} via
\texttt{List.all\_eq\_true}. The reason is that in core Lean the type
\texttt{Fin m} $\times$ \texttt{Fin m} carries no \texttt{Fintype} instance, so
a universally quantified edge condition is not decidable and \texttt{decide}
cannot close it. Second, the inequality $\chic \le \chi$ is formalised
generally, as \texttt{proper\_is\_pq}: a proper $k$-colouring is a
$(k,1)$-colouring, by \texttt{circDist\_pos} and \texttt{circDist\_lt}. The
properness of $c_{61}$ is then \emph{derived} from
\texttt{gp61\_has\_2\_1\_colouring} through \texttt{isPQ\_iff}, not asserted,
and fed back through \texttt{proper\_is\_pq}; both directions of the bridge are
therefore exercised by the build.

Finally, $\chic$ itself is not formalised. Since $\chic(G)$ is the infimum of
the ratios realised by circular colourings, exhibiting a single realised ratio
below $a/b$ already gives $\chic(G) < a/b$; the formalisation states exactly
that constructive form, \texttt{HasColouringBelow}, which keeps the development
inside core Lean. What is \emph{not} formalised is the structural argument of
Section~\ref{sec:structural} --- Theorem~\ref{thm:chi3} and
Corollary~\ref{cor:no4} --- because it rests on Brooks' theorem
(Lemma~\ref{lem:brooks}). Brooks' theorem is neither formalised nor assumed: it
does not occur in the Lean source in any form. This is deliberate, and it does
not weaken the refutation, which is carried by the two certificates alone; but
a reviewer should be aware that the stronger claim ``$\chic \le 3$ for every
admissible $(n,k)$'' remains LaTeX-only.

\section{Reproducibility}

Pure Python~3, standard library only:

\begin{quote}\ttfamily
\$ python3 reproduce.py
\end{quote}

The script (i) builds $\Gn$ for all admissible pairs and verifies
connectedness, $\Delta = 3$, and $\chi \le 3$; (ii) computes exact $\chic$ for
small cases by exhaustive $(p,q)$-search with $p \le |V|$, and prints the
certifying colourings with a violation count of $0$; (iii) sweeps
$5 \le n \le 20$ and reports $96$ cases with $89$ refuted. The structural
conclusion, however, is the proof of Theorem~\ref{thm:main} and rests on
Brooks' theorem, not on the computation.

\end{document}
